\chapter{Direct curved differentiation and surface geometry}
\label{app:geometry-extensions}

This appendix records the mathematics behind two isolated investigations:
direct physical differentiation of a mapped DG function, and first-order
OFDG on a spherical surface. They answer different questions. The first
changes how derivatives of a function on a volume element are evaluated.
The second changes the domain from a volume to a two-dimensional surface
embedded in three-dimensional space. A curved volume element remains a
full-dimensional region; a surface element does not.

\ScopeLimit{The production method remains the recursively projected volume
construction in \cref{sec:curvilinear-mathematics}. Direct differentiation was
investigated in a separate operator experiment. The surface implementation
is an isolated P1/Q1 rotational-transport example, using only first surface
derivatives. Higher-order surface OFDG and general production surface
support are outside the scope of this thesis. The final section explains
that boundary; it does not specify an implemented extension.}

The derivations below introduce their building blocks before combining them.
In particular, a coordinate transformation, a physical $L^2$ projection,
and a geometric projection onto a tangent plane are three distinct operations.
Numerical comparisons remain in the separate supervisor brief; this appendix
does not turn preliminary operator measurements into a PDE accuracy claim.

\section{Direct physical differentiation of the original DG function}
\label{sec:appendix-direct}

\subsection{The function is already defined in physical space}
Let $K\subset\mathbb R^d$ be a volume element and let
$T:\widehat K\longrightarrow K$ be a smooth invertible map. The reference
coordinate is $\xi=(\xi_1,\ldots,\xi_d)$ and the physical coordinate is
$x=(x_1,\ldots,x_d)=T(\xi)$. Write $S=T^{-1}$ for the inverse map.
For coefficients $U_b$ and reference basis functions $\widehat\phi_b$,
\begin{equation}
 \widehat u_h(\xi)=\sum_{b=1}^N U_b\widehat\phi_b(\xi),
 \qquad u_h(x)=\widehat u_h(S(x)).
 \label{eq:app-pullback}
\end{equation}
The second expression is the physical function, not an approximation to a
second representation of it. Its physical basis is simply
$\phi_b(x)=\widehat\phi_b(S(x))$. Thus there is no need to project the input
into a new space merely to regard it as a function of physical position.
The original coefficients already define that function exactly, subject to
floating-point evaluation and the chosen geometry representation.

On a curved element the $\phi_b$ generally are not physical polynomials.
They remain legitimate basis functions. We can evaluate and differentiate
such compositions without requiring their derivatives to lie in the
original finite-dimensional space. What we lose is the ability to store
an arbitrary exact derivative using just another coefficient vector in that
same space. Pointwise evaluation and finite-element representation are
therefore different requirements.

This distinction also separates the present experiment from reconstructing
$u_h$ as a polynomial in $x$. A physical-polynomial $L^2$ reconstruction would
usually change the input function before differentiation. Direct composition
with $S$ does not. All derivatives in this section are element-interior
classical derivatives. We do not differentiate the discontinuous global DG
function in the distributional sense.

\subsection{First derivatives: why the Jacobian is still needed}
Define $J_{ar}=\partial T_a/\partial\xi_r$ and
$A_{ri}=(J^{-1})_{ri}$. Indices $a,i,j$ refer to physical coordinates;
$r,s$ refer to reference coordinates. The inverse function identity
$T(S(x))=x$ gives $DS(x)=J(S(x))^{-1}$. The chain rule then yields
\begin{equation}
 \partial_{x_i}u_h(x)
 =\sum_{r=1}^d \partial_{\xi_r}\widehat u_h(S(x))A_{ri}(S(x)).
 \label{eq:app-first-direct}
\end{equation}
The entries of $A$ describe how the reference position changes when one
physical coordinate changes. Mapping the function to physical space has not
removed this dependence: it has placed it inside the composition in
\cref{eq:app-pullback}. If we differentiate the whole composition
symbolically or by automatic differentiation, the chain rule supplies these
factors automatically. If we differentiate the reference polynomial by hand,
we must supply them explicitly. Omitting them would differentiate a different
function with respect to a different coordinate.

A simple scaling illustrates this. If $x=2\xi$ and
$\widehat u_h(\xi)=\xi$, then $u_h(x)=x/2$. Its physical derivative is $1/2$,
whereas the reference derivative is $1$. The inverse Jacobian is precisely
the missing factor. Curvature adds variation of that factor with position.

\subsection{Second derivatives: differentiate the geometric factor too}
The Hessian of a scalar function is the matrix of all its second partial
derivatives. In two dimensions it is
\begin{equation}
 H_xu_h=
 \begin{pmatrix}
 \partial_{xx}u_h&\partial_{xy}u_h\\
 \partial_{yx}u_h&\partial_{yy}u_h
 \end{pmatrix}.
\end{equation}
This is notation for familiar derivatives, rather than a new differentiation
operation. Applying the product and chain rules to
\cref{eq:app-first-direct} gives, with all reference quantities evaluated at
$\xi=S(x)$,
\begin{align}
 \partial_{x_j}\partial_{x_i}u_h
 &=\sum_{r,s=1}^d A_{sj}A_{ri}
                    \partial_{\xi_s}\partial_{\xi_r}\widehat u_h
   +\sum_{r,s=1}^d A_{sj}(\partial_{\xi_s}A_{ri})
                    \partial_{\xi_r}\widehat u_h.
 \label{eq:app-second-direct}
\end{align}
The first sum differentiates the reference function. The second differentiates
the geometric factor in its first physical derivative. On an affine element
$A$ is constant, so the second sum vanishes. On a curved element it generally
does not. Applying the first-derivative transformation twice while treating
$A$ as constant misses exactly this contribution.

The required inverse-Jacobian derivative follows by differentiating
$JA=I$:
\begin{equation}
 \partial_{\xi_s}A=-A(\partial_{\xi_s}J)A.
\end{equation}
Thus a direct second derivative needs second derivatives of the geometry,
in addition to first and second derivatives of the reference solution.
This requirement is different from the production recursion, which uses
first geometry derivatives to assemble each projected first-derivative
matrix and then differentiates the represented projection again.

\subsection{A one-dimensional example through the third derivative}
Consider
\begin{equation}
 x=T(\xi)=\xi+\varepsilon\xi(1-\xi),\qquad
 0<|\varepsilon|<1,\qquad j(\xi)=T'(\xi)>0.
\end{equation}
Let $f(\xi)=\widehat u_h(\xi)$. Since $d/dx=j^{-1}d/d\xi$,
repeated application of the product rule gives
\begin{align}
 u_h'&=\frac{f'}{j},\\
 u_h''&=\frac{f''}{j^2}-\frac{f'j'}{j^3},\\
 u_h'''&=\frac{f'''}{j^3}
 -\frac{3f''j'}{j^4}-\frac{f'j''}{j^4}
 +\frac{3f'(j')^2}{j^5}.
 \label{eq:app-third-one-dimensional}
\end{align}
Primes on $f$ and $j$ mean reference derivatives; primes on $u_h$ mean
physical derivatives. These equations are evaluated at $x=T(\xi)$.
For $f(\xi)=\xi$, the reference second and third derivatives vanish, but
$j=1+\varepsilon-2\varepsilon\xi$, $j'=-2\varepsilon$, and $j''=0$ give
\begin{equation}
 u_h'=j^{-1},\qquad
 u_h''=2\varepsilon j^{-3},\qquad
 u_h'''=12\varepsilon^2j^{-5}.
\end{equation}
A reference-linear function can therefore have nonzero physical derivatives
of every order. This is not numerical error: it is a property of the
represented function on the curved element. The powers of $j^{-1}$ also
explain why nearly singular maps can amplify derivatives and numerical error.

\subsection{Higher derivatives from a local inverse Taylor series}
\label{sec:appendix-jets}
Writing every product-rule term separately becomes cumbersome in several
dimensions. The isolated experiment instead stores a short multivariate
Taylor polynomial, often called a \emph{Taylor jet}. This polynomial is
local to one evaluation point; it is not a new DG approximation or an
$L^2$ projection of the solution.

Choose a reference point $\xi_0$ and its physical image $x_0=T(\xi_0)$.
Let $h>0$ be a length scale and introduce a dimensionless physical increment
$\eta=(\eta_1,\ldots,\eta_d)$ by $x=x_0+h\eta$. We seek the Taylor
polynomial $q(\eta)$ of the inverse map about this point, through total
degree $m$, satisfying
\begin{equation}
 q(0)=\xi_0,\qquad
 T(q(\eta))=x_0+h\eta+\mathcal O(|\eta|^{m+1}).
 \label{eq:app-inverse-series}
\end{equation}
Substitute it into the original reference function and retain terms through
that degree:
\begin{equation}
 \widehat u_h(q(\eta))
 =\sum_{|\alpha|\le m}c_\alpha\eta^\alpha
    +\mathcal O(|\eta|^{m+1}).
\end{equation}
Here a multi-index is a tuple of nonnegative integers
$\alpha=(\alpha_1,\ldots,\alpha_d)$. Its total order,
factorial and monomial are
\begin{equation}
 |\alpha|=\sum_{i=1}^d\alpha_i,\qquad
 \alpha!=\prod_{i=1}^d\alpha_i!,\qquad
 \eta^\alpha=\prod_{i=1}^d\eta_i^{\alpha_i}.
\end{equation}
For example, $\alpha=(2,1)$ means two $x$ derivatives and one $y$
derivative; $|\alpha|=3$ and $\alpha!=2$. Comparing with the Taylor
expansion in physical coordinates shows that
\begin{equation}
 \partial_x^\alpha u_h(x_0)=\frac{\alpha!}{h^{|\alpha|}}c_\alpha.
 \label{eq:app-jet-coefficient}
\end{equation}
The factorial and length-scale factors are essential. A Taylor coefficient
is not in general equal to the derivative it encodes.

Polynomial addition is coefficientwise. Multiplication multiplies all pairs
of monomials and discards products whose total degree exceeds $m$. Terms of
higher degree cannot contribute to lower-degree coefficients under these
operations. Consequently this truncation preserves the derivatives through
order $m$ at the base point in exact arithmetic. It does not promise an
accurate approximation over a finite neighbourhood of arbitrary size.

For two dimensions and $m=3$, ten coefficients suffice: one constant, two
linear, three quadratic and four cubic coefficients. The experiment computes
the inverse series using the fixed base-point Jacobian $J_0=J(\xi_0)$:
\begin{align}
 q^{(0)}(\eta)&=\xi_0,\\
 q^{(\ell+1)}(\eta)&=
 \left[q^{(\ell)}(\eta)-J_0^{-1}
 \{T(q^{(\ell)}(\eta))-x_0-h\eta\}\right]_{\le m}.
 \label{eq:app-inverse-iteration}
\end{align}
The brackets mean truncation by total degree. To see why this works,
suppose the current inverse-series error begins at degree $r$. Near the
base point, the linear part of the composition error is $J_0$ times that
error. Multiplication by $J_0^{-1}$ cancels its degree-$r$ term; variation
of the Jacobian contributes only higher-degree terms. Each correction thus
fixes at least one additional degree. Three corrections recover degrees
one through three in exact arithmetic; the prototype uses four to include
a redundant residual correction. This is a local series construction, not
a global guarantee of convergence for arbitrary physical displacements.

The geometry fixtures are explicit polynomial maps, so their composition
with a jet is available directly. The reference finite-element polynomial
is expressed in centred monomials by a basis conversion, then evaluated
at $q(\eta)$. That conversion preserves the represented function in exact
arithmetic; it is not a physical-polynomial reconstruction. Its conditioning,
the inverse-series arithmetic and the factor $h^{-|\alpha|}$ can nevertheless
amplify floating-point errors.

\subsection{What direct differentiation improves, and what it cannot fix}
Let $\mathcal R_\alpha u_h$ denote the production recursively projected
derivative in its chosen order, and let $u$ be an underlying smooth exact
field. Then
\begin{align}
 \mathcal R_\alpha u_h-\partial^\alpha u
 &=\underbrace{\mathcal R_\alpha u_h-\partial^\alpha u_h}
                 _{\text{derivative-operator discrepancy}}
   +\underbrace{\partial^\alpha(u_h-u)}
                 _{\text{differentiated input error}},\\
 \partial^\alpha u_h-\partial^\alpha u
 &=\partial^\alpha(u_h-u).
 \label{eq:app-error-split}
\end{align}
Direct differentiation removes the intermediate-projection contribution.
It does not remove error already present in $u_h$, and differentiation can
magnify that error. The two terms in the first line may reinforce or cancel
each other. Removing the first term therefore need not decrease the total
error against $u$, even though it gives a more faithful derivative of $u_h$.
Intermediate projection can also suppress some components of the input error.
These identities explain why an exactly represented test field is useful
for isolating operator error, and why it cannot establish improved PDE
performance on its own.

For sufficiently smooth full-dimensional maps and inputs, exact mixed
Cartesian derivatives commute. The $xxy$, $xyx$ and $yxx$ derivatives of
$u_h$ agree. The jet represents them by the single coefficient with
$\alpha=(2,1)$; it does not independently calculate three differentiation
paths. By contrast, successive finite-space projections can make those
paths differ, as derived in \cref{sec:curved-recursive-derivatives}.
Agreement built into a representation is not an independent correctness
check, which is why the isolated experiment also used physical-coordinate
finite differences with an independent inverse-map calculation.

In principle, derivatives required only at face quadrature points could be
evaluated directly there and their one-sided jumps formed afterwards.
The OFDG damping still uses its physical nested projections; changing the
sensor derivative evaluator would not remove those projections. Such an
integration would require geometry derivatives or a suitable differentiable
map evaluator, pointwise derivative evaluation and verification of the
resulting damping rates. The experiment supports two explicit polynomial
geometry families in two dimensions, through third derivatives. It does not
provide this interface for arbitrary MFEM maps, NURBS or production 3D
filtering, and it has not replaced the current method.

\section{The mathematics needed for the first-order surface example}
\label{sec:appendix-surface}

\subsection{Two surface coordinates and three physical components}
A surface patch $K\subset\Gamma\subset\mathbb R^3$ is parameterized by
$X(\xi_1,\xi_2)$. Its Jacobian has two columns and three rows:
\begin{equation}
 J=(e_1\ e_2),\qquad e_a=\partial_{\xi_a}X\in\mathbb R^3,
 \qquad G=J^{\mathsf T}J.
 \label{eq:app-surface-metric}
\end{equation}
The columns $e_1,e_2$ span the tangent plane. They need not have length one
or be perpendicular. The $2\times2$ matrix $G$, called the metric matrix,
records their lengths and mutual angle: $G_{ab}=e_a\cdot e_b$. For a
regular surface patch these vectors are independent, so $G$ is positive
definite. Although the surface lies in three-dimensional space, it has only
two independent directions of motion. There is no square inverse $J^{-1}$.

A surface DG function is represented as
$u_h(X(\xi))=\widehat u_h(\xi)$. Surface values do not specify how a function
changes when leaving the surface in the normal direction. This is why an
ordinary ambient $x$, $y$ or $z$ derivative is not determined by the surface
function alone. The tangential gradient is determined: it measures changes
along admissible surface directions.

\subsection{Deriving the tangential gradient}
Seek a tangent vector $g$ whose dot product with a tangent displacement
reproduces the directional change of $u$. A reference displacement
$\delta\xi$ induces $\delta x=J\delta\xi$. Therefore
\begin{equation}
 g\cdot J\delta\xi=(\nabla_\xi\widehat u)\cdot\delta\xi
 \quad\hbox{for every }\delta\xi.
\end{equation}
This means $J^{\mathsf T}g=\nabla_\xi\widehat u$. Because $g$ is tangent,
write $g=Jc$ for a two-component coefficient vector $c$. The resulting
system is $Gc=\nabla_\xi\widehat u$, giving
\begin{equation}
 \boxed{\nabla_\Gamma u=J G^{-1}\nabla_\xi\widehat u.}
 \label{eq:app-surface-gradient}
\end{equation}
The result has three Cartesian components, but lies in a two-dimensional
tangent plane. Keeping three components avoids constructing and comparing
a separate orthonormal tangent basis on every element.

An equivalent physical description uses a unit normal $n$ and the matrix
\begin{equation}
 P=I-nn^{\mathsf T}=J G^{-1}J^{\mathsf T}.
\end{equation}
For a vector $a$, $Pa=a-n(n\cdot a)$ removes its normal component. If
$U$ is any smooth ambient extension of $u$, then
\begin{equation}
 \nabla_\Gamma u=P\nabla U.
 \label{eq:app-tangent-projection}
\end{equation}
Different extensions can have different normal derivatives. Their tangential
derivatives agree, because they have identical values along every surface
curve. Thus the right-hand side does not depend on an arbitrary choice of
extension. The matrix $P$ is an exact geometric projection of a vector at
one point; it is not the finite-element $L^2$ projection $\Pi_K$.

\subsection{Surface area, edge length and the co-normal}
A small reference parallelogram has physical area
$|e_1\times e_2|$ times its reference area. Consequently
\begin{equation}
 dS=\sqrt{\det G}\,d\xi_1d\xi_2,\qquad
 |K|=\int_{\widehat K}\sqrt{\det G}\,d\xi.
\end{equation}
A reference edge parameterized by $\widehat\gamma(s)$ has physical tangent
$J\widehat\gamma'(s)$ and line element
\begin{equation}
 d\ell=|J\widehat\gamma'(s)|\,ds.
\end{equation}
A flux leaves a surface cell through its edge. Its relevant outward direction
must be tangent to the surface and perpendicular to the edge. This unit
vector is the \emph{co-normal} $\mu$, distinct from the surface normal $n$.
If $\widehat\nu$ is the outward unit reference-edge normal, then
\begin{equation}
 \mu=\frac{J G^{-1}\widehat\nu}
              {|J G^{-1}\widehat\nu|},\qquad
 \mu\,d\ell=\sqrt{\det G}\,J G^{-1}\widehat\nu\,d\widehat\ell.
 \label{eq:app-conormal}
\end{equation}
Indeed, its dot product with $J\widehat\tau$ is zero whenever
$\widehat\nu\cdot\widehat\tau=0$, so it is perpendicular to the edge
within the tangent plane. The second identity combines this direction with
the physical edge scale. It is useful for conservative flux assembly.
\begin{figure}[tbp]
 \centering
 \begin{tikzpicture}[scale=1.05,>=Stealth]
  \fill[blue!7] (-2.8,-.9)--(.3,-.9)--(1.7,.9)--(-1.4,.9)--cycle;
  \draw[black!40] (-2.8,-.9)--(.3,-.9)--(1.7,.9)--(-1.4,.9)--cycle;
  \draw[very thick,blue!65!black] (.3,-.9)--(1.7,.9);
  \coordinate (p) at (1,0);
  \fill (p) circle (1.8pt);
  \draw[->,thick] (p)--(1,1.65) node[above] {$n$: surface normal};
  \draw[->,thick,orange!80!black] (p)--(3.05,0)
    node[right] {$\mu$: outward co-normal};
  \draw[->,thick,blue!65!black] (p)--(1.58,.75)
    node[right] {$\tau$: edge tangent};
  \node at (-1.1,-.2) {tangent plane};
  \node[below] at (.6,-.75) {cell edge};
 \end{tikzpicture}
 \caption{Local schematic at a surface-cell edge, drawn in perspective.
 The edge tangent $\tau$ and outward co-normal $\mu$ lie in the tangent
 plane and are perpendicular in physical space. The surface normal $n$
 points out of that plane. Transport across the edge uses $v\cdot\mu$,
 not $v\cdot n$.}
 \label{fig:appendix-conormal}
\end{figure}
Figure~\ref{fig:appendix-conormal} distinguishes these three directions.
On the smooth sphere, neighbouring cells have opposite co-normals at their
common edge. For a merely piecewise-planar approximation of a general
surface, this statement requires additional care because the two tangent
planes can differ; the prototype uses an exact smooth radial sphere map.

\subsection{The exact sphere map used in the prototype}
Let $F_h(\xi)$ denote the base element map from an octahedral or cubed-sphere
mesh, with $F_h(\xi)\ne0$. The actual geometry used at quadrature points is
\begin{equation}
 X(\xi)=\frac{F_h(\xi)}{|F_h(\xi)|},\qquad
 J=\frac{1}{|F_h|}(I-XX^{\mathsf T})\,DF_h.
 \label{eq:app-radial-map}
\end{equation}
The second formula follows by differentiating the normalization. Every
mapped point has unit length. The tangent projector is $I-XX^{\mathsf T}$,
since $X$ itself is the outward unit normal on the unit sphere.
The base mesh supplies connectivity and reference coordinates; assembly,
projection, error measurement and visualisation use the radial map.
This is not the same as solving on planar facets whose vertices happen to
lie on the sphere. It also is not a general change to MFEM's geometry API.

The metric, gradient and co-normal constructions above apply to any smooth
regular embedded surface patch. The particular radial map and rotation
problem are sphere-specific. General mathematical formulas do not imply
that the prototype implements all such surfaces or their interfaces.

\subsection{Surface transport and the conservative DG form}
For a tangential velocity $v$, the conservative transport equation is
\begin{equation}
 \partial_t u+\operatorname{div}_\Gamma(vu)=0.
\end{equation}
Surface divergence measures net tangential flux leaving a small surface
region per unit area. For a test function $w_h$, integration by parts on
a cell gives
\begin{equation}
 \int_K w_h\partial_tu_h\,dS
 -\int_K \nabla_\Gamma w_h\cdot v u_h\,dS
 +\int_{\partial K}w_h\widehat f\,d\ell=0.
 \label{eq:app-surface-weak}
\end{equation}
On an edge oriented outward from $K^-$, write $b=v\cdot\mu^-$ and
let $u^-,u^+$ be the two traces. The upwind flux is
\begin{equation}
 \widehat f=b\begin{cases}u^-,&b\ge0,\\u^+,&b<0.\end{cases}
\end{equation}
The same numerical flux enters the neighbouring equation with the opposite
sign. Choosing $w_h=1$ cancels interior fluxes when summing over cells,
which explains global mass conservation on the closed sphere. The normal
speed $v\cdot n$ cannot replace $b$: tangential motion has $v\cdot n=0$
even while it transports material across cell edges.

The example uses $v=(-y,x,0)$, rotation about the $z$ axis. This velocity is
tangent and surface-divergence-free. Tracing a characteristic backwards gives
\begin{equation}
 u(x,y,z,t)=u_0(x\cos t+y\sin t,-x\sin t+y\cos t,z).
\end{equation}
At $T=\pi/2$ the pattern has completed a quarter rotation. This provides an
exact reference without solving a second numerical PDE. The surface is
closed, so there is no physical boundary condition to impose.

\subsection{Projected first derivatives and the order-one damping}
The isolated solution space is discontinuous P1 on triangles or Q1 on
quadrilaterals. Let $\phi_b=\widehat\phi_b\circ X^{-1}$ on a patch, where
$X^{-1}$ means its local inverse from the surface patch to its reference
coordinates. For Cartesian component $a\in\{1,2,3\}$, define
\begin{equation}
 g_{K,a}=\Pi_K[(\nabla_\Gamma u_h)_a].
\end{equation}
Here $\Pi_K$ is the physical surface $L^2$ projection into that same local
scalar DG space. Testing with each basis function gives
\begin{align}
 M^K Z_a&=B_a^K U,\\
 M^K_{ij}&=\int_{\widehat K}\widehat\phi_i\widehat\phi_j
                            \sqrt{\det G}\,d\xi,\\
 (B_a^K)_{ij}&=\int_{\widehat K}\widehat\phi_i
       (J G^{-1}\nabla_\xi\widehat\phi_j)_a\sqrt{\det G}\,d\xi.
\end{align}
The vector $Z_a$ contains coefficients of the projected scalar component.
Although the exact gradient is tangent at every point, separately projecting
its three components need not preserve pointwise tangency: the tangent
plane varies across the element. The prototype uses these three projected
components as its sensor variables. It does not project the resulting
vector tangentially a second time.

For an edge $F$, brackets denote the difference of the two one-sided traces
at the same physical point. The two edge-averaged jump measures are
\begin{equation}
 J_{F,0}=\frac1{|F|}\int_F|[u_h]|\,d\ell,\qquad
 J_{F,1}=\frac1{|F|}\int_F\sum_{a=1}^3|[g_a]|\,d\ell.
\end{equation}
The second expression retains a componentwise absolute sum in fixed ambient
Cartesian axes. It is not a rotationally invariant vector norm, and no such
invariance is claimed. Changing reference coordinates consistently does
not change the exact geometric gradient, but changing the ambient Cartesian
axes can change this particular sensor sum.

Let $h_K=\sqrt{|K|/|\widehat K|}$, let $\beta_F$ be the sampled maximum
of $|v\cdot\mu|$ on the edge, and let $A$ be the sampled global maximum
of $|u_h-\overline u_\Gamma|$, where $\overline u_\Gamma$ is the global
surface mean. Setting $k=1$ in \cref{eq:ofdg-damping-rate} gives
the prefactors $(2l+1)/(2l!)$: $1/2$ for $l=0$ and $3/2$ for $l=1$.
The only nonconstant shell receives both rates. With the surface jump
measures above, the tested rate and update are
\begin{align}
 \lambda_K&=\frac1A\sum_{F\subset\partial K}\beta_F
              \left(\frac{J_{F,0}}{2h_K}+\frac32J_{F,1}\right),\\
 u_h^{\mathrm{new}}&=\overline u_K+
        \exp(-\Delta t\lambda_K)(u_h^{\mathrm{old}}-\overline u_K),
 \qquad \overline u_K=\frac1{|K|}\int_Ku_h\,dS.
 \label{eq:app-surface-damping}
\end{align}
Both terms in the rate have units of inverse time after division by $A$:
one jump is divided by a length, and the derivative jump already has one
inverse length. The near-constant guard sets damping inactive when
$A\le10^{-12}\max(1,\text{sampled volume magnitude})$, avoiding division
by a vanishing amplitude. The factor is applied once after a complete RK4
step, with the rate evaluated from that step's input to the filter.

The cell mean is preserved because the damped difference integrates to
zero. On Q1, all nonconstant modes, including the bilinear reference mode,
belong to this damped tail. Only first derivatives enter this construction;
there is no repeated tangential derivative. Positive surface quadrature
defines the discrete integrals consistently. No positivity limiter or KXRCF
is included, and the mean-preserving update alone does not prove positivity.

\section{Why higher surface derivatives are a separate extension}
\label{sec:appendix-surface-boundary}
\ScopeLimit{This section is explanatory mathematics only. Second- and
higher-order surface derivative sensors have not been implemented or adopted.
It records why increasing the polynomial order would require a new decision,
rather than merely enabling the volume derivative recursion.}

\subsection{Ambient and surface second derivatives answer different questions}
Let $U(x,y,z)=x$ in ambient space and restrict it to the unit sphere:
$u=U|_\Gamma$. Certainly $U_x=1$ and $U_{xx}=0$. Those derivatives vary
$x$ while holding $y,z$ fixed, which generally leaves the sphere. Along
the unit-speed great circle
\begin{equation}
 \gamma(s)=(\cos s,\sin s,0),\qquad
 u(\gamma(s))=\cos s,
\end{equation}
the first and second derivatives with respect to distance $s$ are
$-\sin s$ and $-\cos s$. There is no contradiction. For any smooth ambient
function and path, the ordinary chain rule gives
\begin{equation}
 \frac{d^2}{ds^2}U(\gamma(s))
 =\gamma'(s)^{\mathsf T}(H_U)\gamma'(s)
       +\nabla U\cdot\gamma''(s).
 \label{eq:app-path-second}
\end{equation}
The first term measures change of the gradient; the second accounts for
change in the direction of travel. For $U=x$, the first term is zero and
$\nabla U=(1,0,0)$, while
$\gamma''=(-\cos s,-\sin s,0)$. The second term therefore equals
$-\cos s$. A straight line and a curved path through the same ambient field
need not produce the same second derivative along distance.

\subsection{A physical description of the surface Hessian}
The intrinsic surface Hessian describes how the tangential gradient changes
along the surface, with its output compared in the tangent plane. If
$g=\nabla_\Gamma u$, its ambient matrix representation is
\begin{equation}
 H_\Gamma u=P\,Dg\,P.
 \label{eq:app-surface-hessian}
\end{equation}
Here $Dg$ has entries $(Dg)_{ij}=\partial_{x_j}g_i$ for a smooth local
extension. The right $P$ restricts the differentiation direction to tangent
vectors. The left $P$ keeps only the tangent part of the resulting change.
The value on tangent directions is independent of how $g$ was extended
away from the surface. These geometric projections do not introduce an
$L^2$ approximation.

Since $g=P\nabla U$, differentiating $g$ also differentiates $P$. Surface
curvature appears precisely because the tangent plane changes with position.
Using only $P H_U P$ would omit that effect. On the unit sphere, write
$p=(x,y,z)$, so $P=I-pp^{\mathsf T}$. For $u=x$,
\begin{equation}
 g=e_x-xp,\qquad H_\Gamma u=-xP.
\end{equation}
To verify the latter, take a tangent vector $t$. Differentiating $g$ in that
direction gives $-t_xp-xt$; tangential projection removes the first term
and leaves $-xt$. Thus for a unit tangent $t$,
$t^{\mathsf T}H_\Gamma u\,t=-x$. Along the great circle this is
$-\cos s$, exactly as above. In particular, it is a surface derivative,
not the ambient Cartesian quantity $U_{xx}$.

A general surface curve also has a tangential acceleration. Its second
derivative of $u$ contains the corresponding first-gradient contribution.
For a geodesic, meaning a surface path with zero tangential acceleration,
that contribution vanishes, so the surface Hessian along its tangent gives
the second derivative along the path. Great circles parameterized by arc
length are such paths on the sphere. This qualification matters: not every
curve on a surface measures the Hessian alone.

\subsection{The same construction in reference coordinates}
The reference-coordinate description expresses the same geometric quantity
in the moving basis $e_a=\partial_aX$. It is not another physical derivative
notion. Write $G_{ab}=e_a\cdot e_b$ and let $G^{ab}$ denote entries of
$G^{-1}$, not componentwise reciprocals. Define the coefficients
\begin{equation}
 C^c_{ab}=\sum_{d=1}^2G^{cd}\,e_d\cdot\partial_a\partial_bX.
\end{equation}
These are commonly called Christoffel symbols. Their concrete meaning is
that the tangential part of the changing basis vector $\partial_a e_b$
is $\sum_c C^c_{ab}e_c$. The remaining part points normally to the surface.

The scalar first derivative is $\partial_b\widehat u=g\cdot e_b$.
Differentiating this identity with respect to $\xi_a$ gives
\begin{equation}
 \partial_a\partial_b\widehat u
   =(\partial_a g)\cdot e_b+g\cdot\partial_a e_b.
\end{equation}
The second term measures change of the basis vector itself. Subtracting
its tangential contribution isolates change of the gradient:
\begin{equation}
 H_{ab}=\partial_a\partial_b\widehat u
            -\sum_{c=1}^2 C^c_{ab}\partial_c\widehat u.
 \label{eq:app-coordinate-hessian}
\end{equation}
These are the surface Hessian evaluated on the coordinate tangent vectors.
Recovering its ambient matrix representation uses
$H_\Gamma u=J G^{-1}H G^{-1}J^{\mathsf T}$. Hence the coordinate formula
and \cref{eq:app-surface-hessian} describe the same object. The correction
is required even though the reference second derivatives of the scalar
function can be evaluated easily.

\subsection{Why third order is not merely another scalar multi-index}
Differentiating a scalar once gives a gradient; differentiating that gradient
intrinsically gives a two-slot object, the Hessian. A further derivative must
account for changes in both of its tangent arguments. With the first index
specifying the differentiation direction, its coordinate components are
\begin{equation}
 T_{abc}=\partial_aH_{bc}
      -\sum_{d=1}^2 C^d_{ab}H_{dc}
      -\sum_{d=1}^2 C^d_{ac}H_{bd}.
 \label{eq:app-third-covariant}
\end{equation}
Each subtraction corrects for one moving tangent-basis vector. In general,
a tensor with $r$ such arguments acquires $r$ corresponding correction
terms on differentiation. Direct third-order evaluation generally requires
geometry derivatives through third order. Repeatedly differentiating three
Cartesian components as if they were unrelated scalar fields does not, by
itself, specify these intrinsic tensor derivatives.

There is also a fundamental difference from Cartesian volume derivatives.
The surface Hessian of a smooth scalar field is symmetric, but its covariant
third derivative need not be fully symmetric on a curved surface. On the
unit sphere the preceding example gives $H_{bc}=-xG_{bc}$. The intrinsic
derivative preserves the metric, so
\begin{equation}
 T_{abc}=-(\partial_a x)G_{bc}.
\end{equation}
At the north pole choose orthonormal tangent directions $e_1=e_x$ and
$e_2=e_y$. Then $T_{122}=-1$ whereas $T_{212}=0$. This difference exists
in exact geometry and exact arithmetic. It is not the projection-order
error measured for volume derivatives. Symmetrizing all slots would define
a different sensor quantity, rather than correcting that difference as a
numerical defect.

